\newpage
\chapter{講師登録・講義開講申請}

\section{講師登録申請を受けたとき}

担当分野について他の学生へ伝えられる知識や経験の裏付けは確認するが、競技成績、経験年数又は実績欄の記載を一律の要件にしない。
担当分野が狭い場合は広げる選択肢を助言できるが、それを認可の条件にしない。

\section{講義開講申請を受けたとき}

担当者は、申請者が講師であること、開講申請の必要事項、講義の案内及び受講条件の有無を確認する。
必要事項が揃った申請は、講師の活動を実現するため原則として認可に向けて扱う。
学部や講義区分は、学園規則第７章第５条により講義を探す目安である。
横断する講義は主たる区分を選べばよく、区分の選び直しは助言として伝える。
題材が珍しいこと、範囲が狭いこと又は既存の講義と重なることだけを不認可の理由にしない。
講義回数は講師の予定であり、内容、構成及び進め方は講師の裁量に属する（第６章第３条）。
実施が難しい回は第７章第９条の補講期間の扱いを案内する。

\section{補正、認可及び結果の通知}

記載不足や事実関係の疑義があれば、担当者は申請者に確認し、補正できる事項を具体的に伝える。
学園規則第７章第１０条の中止・変更の事由は開講申請の認可基準を直接定めるものではなく、それだけを申請不認可の根拠として機械的に用いない。
補正後の資料と判断の根拠を教務主事に共有し、講師登録は第６章第２条第２項、開講は第７章第６条第２項に従い、いずれも教務主事全員一致で認可する。
学園長の単独決定で認可に代えない。

申請者には結果を伝える。不認可とする場合は、確認した事実と理由を示す（第５章第６条第２項）。
担当者は申請、補正依頼、各教務主事の判断、結論及び連絡を記録する。
理由の再確認や再審査を求められたときは教務主事間で検討する。
